% !TeX root = ../../Book3.tex
\chapter{Dirichlet \texorpdfstring{\(L\)}{L}-函数与算术级数}
\label{b3:app:D}
\begin{chapterabstract}
按剩余类筛选整数，需要把乘法结构与有限 Fourier 分解结合起来。Dirichlet 特征完成筛选，带特征的 Euler 乘积则把素数的剩余类信息编码为全纯函数。这里从正交关系走到 \(s=1\) 非消失，再证明每个互素算术级数含有无穷多个素数；原始特征的函数方程另由 Gauss 和与 theta 变换推导。
\end{chapterabstract}
\begin{learninggoals}
\begin{itemize}
\item\label{b3:item:D-goal-1}用特征正交关系提取指定的互素剩余类。
\item\label{b3:item:D-goal-2}区分主特征、原始特征与诱导特征。
\item\label{b3:item:D-goal-3}证明非主特征在1点的非消失。
\item\label{b3:item:D-goal-4}追踪函数方程中的奇偶性、模数与 Gauss 和。
\end{itemize}
\end{learninggoals}
前置为第二十五至二十七章。除有限维复线性代数外，不要求一般代数数论。本文中的算术级数定理指素数的无穷性；其带主项的计数渐近还需要整条边界上的非消失分析，这里不把两种结论混称。

\section{有限乘法群的 Fourier 分解}
\label{b3:sec:D01}

设 \(q\ge2\) 为整数，\(G=(\mathbb Z/q\mathbb Z)^\times\)。从 \(G\) 到单位圆的群同态称为群的特征；将它周期延拓到整数，并在 \((n,q)>1\) 时置零，就得到模 \(q\) 的\term[dirichlet-tezheng]{Dirichlet 特征}[Dirichlet character] \(\chi\)。在 \(G\) 上恒等于一的称为主特征 \(\chi_0\)。所有特征都完全乘法，非零值是有限阶单位根。

\begin{proposition}[正交关系]\label{b3:prop:D-orthogonality}
设 \(q\ge1\) 为整数，\(G=(\mathbb Z/q\mathbb Z)^\times\)，特征指 \(G\) 的复单位圆值群同态，在与 \(q\) 不互素的整数上延为0。则模 \(q\) 恰有 \(\varphi(q)\) 个特征。若整数 \(a\) 满足 \((a,q)=1\)，则对每个整数 \(n\)，
\[
 \mathbf1_{n\equiv a\ (q)}=
 \frac1{\varphi(q)}\sum_{\chi\bmod q}\overline{\chi(a)}\chi(n).
\]
并且 \(\sum_{a\in G}\chi(a)\overline{\psi(a)}\) 在 \(\chi=\psi\) 时为 \(|G|\)，否则为零。
\end{proposition}
\begin{proof}
在复向量空间 \(\mathbb C^G\) 上，平移 \(U_bv(a)=v(ab)\) 是彼此交换的酉算子。先对一个算子作酉对角化，其特征子空间被其余算子保持；在这些子空间上继续，即得共同特征向量的正交基。若 \(U_bv=\chi(b)v\)，则 \(v(b)=\chi(b)v(1)\)，非零向量必有 \(v(1)\ne0\)，共同特征子空间是一维，且 \(\chi(bc)=\chi(b)\chi(c)\)。故恰有 \(|G|\) 个互异特征。不同特征正交也可直接验证：选 \(b\) 使 \(\chi(b)\overline{\psi(b)}\ne1\)，平移求和后该和等于自身乘此常数，只能为零。把特征向量除以 \(\sqrt{|G|}\) 后组成酉矩阵，其列也正交，于是 \(\sum_\chi\overline{\chi(a)}\chi(n)=|G|\mathbf1_{a=n}\)。非单位 \(n\) 时两边均为零。
\end{proof}
模4有两个特征，非主特征 \(\chi_4\) 在1、3上分别取1、\(-1\)。模5的单位群由2生成，可分别令 \(\chi(2)=1,\mathrm{i},-1,-\mathrm{i}\)；两个非实特征互为共轭。这些共轭对会在非消失证明中发挥作用。
\ProseParagraphBreak
设 \(q\ge2\) 为整数、\(d\) 为 \(q\) 的正因子；模一的特征在所有整数上恒取一。模 \(d\) 特征可沿单位群的约化映射拉回，再在非单位上置零，称为诱导。不能由真因子模数诱导的特征称为\term[yuanshi-tezheng]{原始特征}[primitive character]。最小诱导模数称导子。例如模8上仅由奇数模4的余数决定的特征，其导子为4；它在模8意义下非原始。原始性负责下面的有限 Fourier 变换不丢失频率。

\ProseParagraphBreak
筛选公式是有限乘法群上的 Fourier 反演，不能换成整数加法群上的同一个求和。以模五为例，取生成元二并令 \(\chi_j(2)=\mathrm{i}^j\)、\(0\le j<4\)。在目标剩余类二处，若 \(n\equiv2^l\pmod5\)，筛选和为
\[
 \frac14\sum_{j=0}^3\overline{\chi_j(2)}\chi_j(n)
 =\frac14\sum_{j=0}^3\mathrm{i}^{j(l-1)}.
\]
当 \(l=1\) 值为一，否则四项几何和为零。若五整除 \(n\)，所有特征值为零，筛选也为零。这说明特征不是只取 \(\pm1\) 的符号函数，非实单位根同样是精确挑出剩余类所需的频率。
\ProseParagraphBreak
一个非主特征的周期和为零也可以不用完整正交矩阵来证明。选单位 \(b\) 使 \(\chi(b)\ne1\)，乘法 \(a\mapsto ba\) 置换整个周期的单位类，故 \(\sum_a\chi(a)=\chi(b)\sum_a\chi(a)\)，只能为零。把 \(\lfloor x/q\rfloor\) 个整周期消去后，\(\sum_{n\le x}\chi(n)\) 只剩不足 \(q\) 个模长至多一的项，所以绝对值不超过 \(q\)。于是部分求和尾界具体为
\[
 \left|\sum_{n>N}\frac{\chi(n)}{n^s}\right|
 \le Cq N^{-\sigma}\left(1+\frac{|s|}{\sigma}\right),\qquad\sigma=\operatorname{Re}s>0.
\]
在实部为正的紧集上，这个界一致趋零；它既说明全纯性，也允许在 \(s=1\) 把 \(L\) 看作通常收敛的级数。


\section{解析延拓与函数方程}
\label{b3:sec:D02}

对 \(\operatorname{Re}s>1\)，定义
\[
 L(s,\chi)=\sum_{n\ge1}\frac{\chi(n)}{n^s}
     =\prod_p(1-\chi(p)p^{-s})^{-1}.
\]
第二十五章的绝对收敛乘法证明逐字适用。若 \(\chi\ne\chi_0\)，一个周期内的和为零，故 \(\sum_{n\le x}\chi(n)=O_q(1)\)，部分求和已把级数全纯延拓到 \(\operatorname{Re}s>0\)。主特征则满足
\[
 L(s,\chi_0)=\zeta(s)\prod_{p\mid q}(1-p^{-s}),
\]
在1点有留数 \(\varphi(q)/q\) 的单极点。

\begin{proposition}[非主特征的整延拓]\label{b3:prop:D-entire-continuation}
设 \(q\ge1\) 为整数，\(\chi\) 为模 \(q\) 的 Dirichlet 特征，\(\chi_0\) 为主特征，\(L(s,\chi)=\sum_{n\ge1}\chi(n)n^{-s}\) 最初定义于 \(\operatorname{Re}s>1\)。若 \(\chi\ne\chi_0\)，则 \(L(s,\chi)\) 延拓为整函数。
\end{proposition}
\begin{proof}
令
\[
 H_\chi(t)=\sum_{n\ge1}\chi(n)e^{-nt}
   =\frac{\sum_{a=1}^q\chi(a)e^{-at}}{1-e^{-qt}}.
\]
分子在0点消失，分母有单零点，故 \(H_\chi\) 在0点全纯。对 \(\operatorname{Re}s>1\)，逐项积分给
\(\Gamma(s)L(s,\chi)=\int_0^\infty H_\chi(t)t^{s-1}\,\mathrm{d}t\)。在 \([0,1]\) 中减去任意多项 Taylor 多项式，积分余项延拓到任意左半平面，而减去的第 \(j\) 项只产生 \(c_j/(s+j)\)。故右侧至多在非正整数有单极点。\(1/\Gamma(s)\) 恰在这些点有单零点，乘后全部可去。在无穷远的积分因 \(H_\chi(t)=O(e^{-t})\) 对一切参数局部一致收敛，因此没有其他奇点。
\end{proof}
若 \(\chi\) 由原始 \(\chi^*\) 诱导，则
\(L(s,\chi)=L(s,\chi^*)\prod_{p\mid q}(1-\chi^*(p)p^{-s})\)。导子中已出现的素数其因子等于一。这些额外因子是讨论零点时必须保留的资料。

\begin{lemma}[原始 Gauss 和]\label{b3:lem:D-gauss-sum}
设 \(q\ge1\) 为整数，\(\chi\) 为模 \(q\) 的原始 Dirichlet 特征，\(n\in\mathbb Z\)，\(\tau(\chi)=\sum_{a=1}^q\chi(a)e^{2\pi\mathrm{i}a/q}\)。则
\[
 \sum_{a=1}^q\chi(a)e^{2\pi\mathrm{i}na/q}=\overline{\chi(n)}\tau(\chi),
 \qquad |\tau(\chi)|=\sqrt q.
\]
\end{lemma}
\begin{proof}
若 \((n,q)=1\)，将 \(a\) 换为 \(n^{-1}a\) 即得第一式。若 \(d=(n,q)>1\)，指数项只依赖 \(a\bmod q/d\)。单位群向模 \(q/d\) 单位群的映射满射：给定与 \(q/d\) 互素的剩余类，逐个避开其余素因子，再用中国剩余定理即可取到与 \(q\) 互素的代表。原始性说明 \(\chi\) 在该映射的核上不恒等于一，否则会下降到真因子模数。选核中 \(b\) 使 \(\chi(b)\ne1\)，把 \(a\) 换为 \(ba\)，指数不变而总和乘以 \(\chi(b)\)，故和为零。最后对 \(n\bmod q\) 作有限 Fourier 的 Parseval 等式，左侧平方和为 \(q\sum_a|\chi(a)|^2=q\varphi(q)\)，右侧为 \(\varphi(q)|\tau(\chi)|^2\)，得到模长。
\end{proof}
令 \(\chi(-1)=(-1)^\varepsilon\)，\(\varepsilon\in\{0,1\}\)，并设 \(\chi\) 原始非主。置
\[
 \theta_\chi(t)=\sum_{n\in\mathbb Z}n^\varepsilon\chi(n)e^{-\pi n^2t/q},
 \quad \Lambda(s,\chi)=\left(\frac q\pi\right)^{(s+\varepsilon)/2}
 \Gamma\left(\frac{s+\varepsilon}2\right)L(s,\chi).
\]
在 \(\varepsilon=0\) 时也令零项为零。把求和按剩余类展开，令 \(f_t(x)=x^\varepsilon e^{-\pi t x^2/q}\)，则
\(\theta_\chi(t)=\sum_{a=1}^q\chi(a)\sum_{k\in\mathbb Z}f_t(a+qk)\)。对步长为 \(q\) 的周期化作 Fourier 展开，系数是 \(q^{-1}\widehat f_t(m/q)\)，所以
\[
 \theta_\chi(t)=\frac1q\sum_{m\in\mathbb Z}\widehat f_t(m/q)
       \sum_{a=1}^q\chi(a)e^{2\pi\mathrm{i}ma/q}
 =\frac{\tau(\chi)}q\sum_{m\in\mathbb Z}\overline{\chi(m)}\widehat f_t(m/q).
\]
Gaussian 及其导数快速衰减，每一个展开都绝对收敛，有限剩余类和也可自由交换。偶情形有
\(\widehat f_t(\xi)=\sqrt{q/t}\,e^{-\pi q\xi^2/t}\)。奇情形利用 \(\widehat{xf}(\xi)=\mathrm{i}\widehat f'(\xi)/(2\pi)\)，得到下面的线性因子。代入频率 \(m/q\)，恰好分别出现整数幂 \(m^0\)、\(m^1\)，因此得到
\[
 \theta_\chi(t)=\frac{\tau(\chi)}{\mathrm{i}^{\varepsilon}\sqrt q}
 t^{-1/2-\varepsilon}\theta_{\bar\chi}(1/t).
\]
这里奇情形的因子来自 \(\widehat{x e^{-\pi tx^2/q}}(\xi)=-(\mathrm{i}q\xi/t)\sqrt{q/t}\,e^{-\pi q\xi^2/t}\)；余数求和把频率变成整数除以 \(q\)，正好留下 \(\mathrm{i}^{-1}t^{-1}\)。偶情形没有这一因子。
\ProseParagraphBreak
两端的 theta 函数在无穷远指数衰减，变换式使其在零点也按 \(e^{-c/t}\) 衰减。因此 Mellin 积分对所有 \(s\) 收敛，且
\[
 2\Lambda(s,\chi)=\int_0^\infty\theta_\chi(t)t^{(s+\varepsilon)/2-1}\,\mathrm{d}t.
\]
这里 \(n\mapsto-n\) 时，\(n^\varepsilon\chi(n)\) 不变，正负项成对贡献，才得到左边的因子二。对正 \(n\)，换元 \(v=\pi n^2t/q\) 后，积分给
\(n^\varepsilon(q/\pi n^2)^{(s+\varepsilon)/2}\Gamma((s+\varepsilon)/2)\)，指标的 \(n\) 次幂正好是 \(n^{-s}\)。这一识别先在 \(\operatorname{Re}s>1\) 完成，再以整函数延拓。
\ProseParagraphBreak
代入 theta 变换并令 \(u=1/t\)，积分中的幂次变为
\[
 t^{(s+\varepsilon)/2-1}t^{-1/2-\varepsilon}\,\mathrm{d}t
 \longmapsto u^{(1-s+\varepsilon)/2-1}\,\mathrm{d}u.
\]
端点反转消去微分的负号，指数衰减保证这次全区间换元对所有 \(s\) 有效。因此完整函数方程为
\begin{equation}\label{b3:eq:D-functional-equation}
 \Lambda(s,\chi)=\frac{\tau(\chi)}{\mathrm{i}^{\varepsilon}\sqrt q}
                      \Lambda(1-s,\bar\chi).
\end{equation}
右侧常数模长为一；共轭特征一般不同于原特征，所以不能无条件写成关于实轴的对称。与不带完成因子的传统写法的对应可核对 NIST DLMF §25.15 的公式(25.15.5)\cite{NationalInstituteofStandardsNIST}。

\section{\texorpdfstring{\(s=1\)}{s=1} 处的非消失}
\label{b3:sec:D03}

实特征的非消失需要正系数级数的一个原则。下面给出证明，以免把核心算术困难藏在“Euler 乘积非零”中：乘积的初始区域并不包含1点。

\begin{lemma}[Landau 的正系数原则]\label{b3:lem:D-landau}
设 \(a_n\in[0,\infty)\ (n\ge1)\)，Dirichlet 级数 \(F(s)=\sum_{n\ge1}a_n n^{-s}\) 的收敛横坐标 \(\sigma_c\in\mathbb R\) 有限，其中 \(s\in\mathbb C\)、\(n^{-s}=e^{-s\log n}\) 用实对数，则 \(F\) 不可能在实点 \(\sigma_c\) 的邻域全纯。
\end{lemma}
\begin{proof}
反设 \(F\) 全纯延拓到 \(B(\sigma_c,\rho)\)。取 \(0<h<\rho/4\)、\(s_0=\sigma_c+h\)、\(R=\rho/2\)。圆盘 \(B(s_0,R)\) 中实部大于 \(\sigma_c\) 的部分在原收敛半平面内，其余部分与 \(\sigma_c\) 的距离小于 \(R+h<\rho\)，也在延拓圆盘中。因此整个 \(B(s_0,R)\) 的函数全纯，且 \(R>h=s_0-\sigma_c\)，Taylor 半径确实跨过实轴收敛边界。非负系数使条件收敛与绝对收敛在实轴相同；在 \(s_0\) 内部可以吸收对数因子后逐项求导。对 \(0<t<R\)，Taylor 展开及逐项微分给
\[
 F(s_0-t)=\sum_{k\ge0}\frac{t^k}{k!}\sum_{n\ge1}a_n(\log n)^k n^{-s_0}
        =\sum_{n\ge1}a_n n^{-s_0}e^{t\log n}.
\]
各项非负，交换双和由单调收敛成立。选 \(s_0-\sigma_c<t<R\)，右侧说明级数在 \(s_0-t<\sigma_c\) 收敛，矛盾。
\end{proof}

\begin{theorem}[Dirichlet 非消失]\label{b3:thm:D-nonvanishing}
设 \(q\ge2\) 为整数，\(\chi\) 为模 \(q\) 的非主 Dirichlet 特征，以 \(L(s,\chi)=\sum_{n\ge1}\chi(n)n^{-s}\)（\(\operatorname{Re}s>0\)）定义其全纯 \(L\)-函数。则 \(L(1,\chi)\ne0\)。
\end{theorem}
\begin{proof}
先设 \(\chi\) 实且非主。\(F(s)=\zeta(s)L(s,\chi)\) 的系数为 \(a_n=\sum_{d\mid n}\chi(d)\)。在素数幂处，若 \(\chi(p)=1,0,-1\)，该系数分别为 \(r+1,1,\mathbf1_{r\text{ 偶}}\)，故全部非负，且 \(a_{m^2}\ge1\)。于是收敛横坐标在 \([1/2,1]\) 内。若 \(L(1,\chi)=0\)，它消去 \(\zeta\) 的唯一极点，\(F\) 在 \(\operatorname{Re}s>0\) 全纯，与 Landau 引理矛盾。
对于所有模 \(q\) 特征，实 \(\sigma>1\) 时有
\[
 \log\prod_\chi L(\sigma,\chi)
 =\sum_{p\nmid q}\sum_{r\ge1}\frac{\sum_\chi\chi(p)^r}{r p^{r\sigma}}\ge0.
\]
内和由正交性等于0或 \(\varphi(q)\)，所以乘积至少为一。若某个非实特征在1点为零，其共轭也在1点为零。主特征至多贡献一个单极点，所有实非主特征已证非零，其余因子全纯；故整个乘积趋于零，与至少为一矛盾。
\end{proof}
实特征与非实特征采用不同机制：前者靠整数平方保证收敛边界确实位于正半平面，后者靠共轭成对的零点压倒唯一极点。若漏掉实情形，乘积论证无法排除一个简单零点恰好抵消主特征的极点。

\ProseParagraphBreak
实情形中起作用的正系数还可以逐素数核验。若 \(a=\mathbf1*\chi\)，则 \(a\) 可乘；\(\chi(p)=-1\) 时 \(a(p^r)=1-1+\cdots+(-1)^r\)，偶指数为一、奇指数为零。因此对平方 \(m^2\)，每一个局部因子都至少一，\(a(m^2)\ge1\)。实 \(\sigma\le1/2\) 时，仅平方项的和便不小于 \(\sum_m m^{-2\sigma}\)，所以发散；而 \(\sigma>1\) 时原双级数绝对收敛，故 \(1/2\le\sigma_c\le1\)。若在一处零点消去了唯一极点，\(\zeta L\) 已在整个正半平面全纯，特别在这个实 \(\sigma_c\) 点全纯，与正系数原则正面矛盾。证明没有要求先知道任何素数在剩余类中的密度。
\ProseParagraphBreak
还可在模三的非主特征上直接看见 \(L(1)\) 的正值。令 \(\chi_3(1)=1,\chi_3(2)=-1\)，则
\[
 L(1,\chi_3)=\sum_{n\ge0}\left(\frac1{3n+1}-\frac1{3n+2}\right)
 =\int_0^1\frac{1-x}{1-x^3}\,\mathrm{d}x
 =\frac\pi{3\sqrt3}.
\]
每对差可写成 \(\int_0^1x^{3n}(1-x)\,\mathrm{d}x\)，非负性允许单调收敛换序；约去 \(1-x\) 后积分为 \(\int_0^1\mathrm{d}x/(1+x+x^2)\)，配方并计算反正切得到最后常数。一般的非消失证明不能靠这种个别积分，但这个算例同时检验了特征取值、分组与边界值的普通收敛意义。
\ProseParagraphBreak
函数方程现在还确定原始非主特征的平凡零点。对 \(s<0\)，\(1-s>1\)，Euler 乘积使右边完成函数非零；左边 Gamma 因子仅在 \(s=-\varepsilon-2m\) 有单极点。因此这些负整数处 \(L\) 必有恰好单零点，其他负实点没有零点。偶特征在零点也有单零，因为函数方程将完成值送到一处，而刚证的 \(L(1,\bar\chi)\ne0\) 使该值非零。诱导特征还带有限 Euler 因子，它们可能添加其他零点，所以这里的结论不能不加原始性就照搬。


\section{算术级数中的素数}
\label{b3:sec:D04}

\begin{theorem}[Dirichlet 算术级数定理]\label{b3:thm:D-dirichlet}
设 \(q\ge1\)、\(a\) 为整数。若 \((a,q)=1\)，则有无穷多个素数 \(p\equiv a\pmod q\)。更精确地，当实 \(\sigma\downarrow1\) 时，
\[
 \sum_{p\equiv a\ (q)}p^{-\sigma}
 =\frac1{\varphi(q)}\log\frac1{\sigma-1}+O_q(1).
\]
\end{theorem}
\begin{proof}
在绝对收敛区定义 Euler 对数，则
\(\log L(\sigma,\chi)=\sum_p\chi(p)p^{-\sigma}+O(1)\)，误差由所有 \(r\ge2\) 的和控制，对 \(\sigma\ge1\) 一致有界。非主特征在1点全纯非零，在其小邻域可取全纯对数；沿实区间与 Euler 对数相差一个固定的 \(2\pi\mathrm{i}\) 整数倍，故有界。主特征则有
\(\log L(\sigma,\chi_0)=\log(1/(\sigma-1))+O_q(1)\)。乘以 \(\overline{\chi(a)}\) 后求和，正交关系选出指定剩余类，得到公式。右侧趋于无穷，左侧不可能只是有限多个素数的和。
\end{proof}
公式描述的是带权和在 \(\sigma=1\) 附近的发散。要推出 \(\pi(x;q,a)\sim x/(\varphi(q)\log x)\)，需要沿整条 \(\operatorname{Re}s=1\) 排除相应奇点，再应用第二十七章的 Tauber 方法。仅凭本节在单点的估计，不能完成这一步。模数不互素的剩余类则完全不同：若 \(d=(a,q)>1\)，该类的任何素数必须整除 \(d\)，因而至多有限个。
\ProseParagraphBreak
继续阅读可从 Apostol《Introduction to Analytic Number Theory》的特征与 \(L\)-级数章节\cite{Apostol1976Introduction}入手；统一的零点估计和算术级数中的素数定理见 Davenport《Multiplicative Number Theory》\cite{Davenport1980Multiplicative}。

\begin{chaptersummary}
特征把剩余类指标分解为乘法函数，Euler 乘积把乘法结构交给复分析。1点非消失是从表示公式走向素数无穷性的关键，而函数方程进一步把右半平面的信息传到左半平面；二者承担不同任务。
\end{chaptersummary}
\BookExercises
\begin{exercise}\label{b3:ex:D-1}
写出模5的四个特征在1、2、3、4上的值。
\end{exercise}
\begin{solution}
按 \(1,2,4,3\) 排列，它们依次为 \((1,1,1,1)\)、\((1,\mathrm{i},-1,-\mathrm{i})\)、\((1,-1,1,-1)\)、\((1,-\mathrm{i},-1,\mathrm{i})\)。换回题目顺序即可。非单位余数0上均为零。
\end{solution}
\begin{exercise}\label{b3:ex:D-2}
计算 \(L(1,\chi_4)\)。
\end{exercise}
\begin{solution}
在 \(0\le r<1\) 上积分几何级数 \(1/(1+x^2)=\sum(-1)^nx^{2n}\)，得 \(\arctan r=\sum(-1)^nr^{2n+1}/(2n+1)\)。令 \(r\uparrow1\)，交错级数的尾项一致受首个未取项控制，故极限为 \(1-1/3+1/5-\cdots=\pi/4\)。
\end{solution}
\begin{exercise}\label{b3:ex:D-3}
模8上由 \(\chi_4\) 诱导的特征是否与 \(\chi_4\) 有不同的 \(L\)-函数？
\end{exercise}
\begin{solution}
没有。两个函数在奇数上相同，在偶数上都为零，Dirichlet 级数完全相同。附加模数只增加素因子2，而 \(\chi_4(2)=0\)，所补 Euler 因子为一。
\end{solution}
\begin{exercise}\label{b3:ex:D-4}
证明 \(\sum_{p\equiv a\ (q)}1/p\) 发散。
\end{exercise}
\begin{solution}
若该正项级数收敛，则对所有 \(\sigma>1\)，\(p^{-\sigma}\le p^{-1}\)，其和一致有界；这与定理~\ref{b3:thm:D-dirichlet}的对数发散矛盾。
\end{solution}
\begin{exercise}\label{b3:ex:D-5}
求 \(\tau(\chi_4)\) 并确定其函数方程的根数。
\end{exercise}
\begin{solution}
\(\tau(\chi_4)=e^{\pi\mathrm{i}/2}-e^{3\pi\mathrm{i}/2}=2\mathrm{i}\)。它是奇特征，\(\varepsilon=1,q=4\)，故根数 \(\tau/(\mathrm{i}\sqrt q)=1\)。于是完成函数满足 \(\Lambda(s,\chi_4)=\Lambda(1-s,\chi_4)\)。
\end{solution}
\begin{exercise}\label{b3:ex:D-6}
证明非主特征的 \(L(0,\chi)=-q^{-1}\sum_{a=1}^q a\chi(a)\)。
\end{exercise}
\begin{solution}
前述生成函数在0点的常数项为 \(H_\chi(0)=-q^{-1}\sum a\chi(a)\)。其 Mellin 变换在 \(s=0\) 的留数等于该常数，而 \(\Gamma(s)\) 在0点留数为一；比较 \(\Gamma(s)L(s,\chi)\) 的留数即得公式。
\end{solution}
